\chapter{Logical Dependencies and Consistency Certificates}

\section{Causal chain of the certificates: finite states \(\to\) restrictions \(\to\) survival \(\to\) infinite section \(\to\) realizations}

The logical chain of the integration does not begin from an Archimedean conclusion.
It is organized in the following order:
\[
\begin{tikzpicture}[baseline=(current bounding box.center),>=Latex,
 node distance=7mm and 8mm,
 n/.style={rounded corners,draw=hmtblue,fill=hmtlightblue,
 minimum height=8mm,align=center,font=\small}]
 \node[n] (finite) {estados finitos\\tipados};
 \node[n,right=of finite] (maps) {restricciones\\compatibles};
 \node[n,right=of maps] (surv) {supervivencia\\y poda};
 \node[n,right=of surv] (lim) {sección\\infinita};
 \node[n,right=of lim] (read) {lectores\\y realizaciones};
 \draw[->,thick] (finite)--(maps);
 \draw[->,thick] (maps)--(surv);
 \draw[->,thick] (surv)--(lim);
 \draw[->,thick] (lim)--(read);
\end{tikzpicture}
\]
A digit certificate verifies a restriction of the last arrow; it does not
retrospectively create the states or replace the existence theorem
for the limit.

\section{Exact arithmetic identities of finite certificates}

The censuses accompanying the nonadic filtration certificate satisfy
\[
 396\cdot6=2376,
 \qquad
 43\cdot9=387,
 \qquad
 432+4\cdot144=1008.
\]
The first equality relates elementary boundary events and trits
emitted per channel; the second counts the phase pairs \((K,K+9)\); the
third reconstructs the total multiplicity of the five regions of
\(\pi\).

The two regional roles that cannot be permuted are
\[
\begin{array}{rcl}
U_\pi^\perp&=&(501,614,498,169,272,272),\\
E_{\rm sig}&=&555555,\qquad \text{multiplicity }432,\\[.35em]
U_\pi^{--}&=&(870,923,810,418,674,521),\\
E_{\rm sig}&=&933717,\qquad \text{multiplicity }144.
\end{array}
\]

\section{Existence and uniqueness of the infinite branch in a finitely ramified tree of surviving prefixes}

\begin{lemma}[Finitely branching tree]
Let \(T\) be a tree with nonempty finite levels such that every
surviving node has a surviving descendant. Then there is an infinite
branch. If, moreover, each node of the selected subtower has exactly
one compatible descendant, the branch is unique.
\end{lemma}

\begin{proof}
Existence is König's lemma applied to the finitely branching tree.
For uniqueness, two distinct branches would have a first level at which they
diverge, yielding two compatible descendants of the same prefix, contrary to the
hypothesis.
\end{proof}

\begin{controlbox}
An execution to depth \(N\) proves only what occurs in the
truncated tree. The infinite claim requires the uniform law that preserves
nonemptiness and compatibility at every level. The main text expressly separates
the two statuses and, for the nonadic filtration, retains its certified uniform
selector.
\end{controlbox}

\section{Closure and Full Power}

For \(A\subseteq X=\varprojlim X_n\), it has been proved
\[
 \overline A=\bigcap_n p_n^{-1}(p_n(A)).
\]
The set of eventually null binary sequences is dense and proper in
\(2^{\mathbb N}\), yet it has all finite shadows. Therefore any
algorithm that receives only \((p_n(A))_n\) cannot distinguish \(A\) from its
closure. This example blocks the false identity
\[
 \varprojlim\mathcal P(X_n)=\mathcal P(\varprojlim X_n).
\]

\section{Typed compatibility of operators}

The difference between the unital inclusion and the marked corner is verified directly:
\[
\begin{aligned}
 \tau_{d+1}(a\otimes I_9)
 &=\frac{\operatorname{Tr}(a)\operatorname{Tr}(I_9)}{9^{d+1}}
 =\tau_d(a),\\
 \tau_{d+1}(a\otimes p_j)
 &=\frac{\operatorname{Tr}(a)}{9^{d+1}}
 =\frac1{9}\tau_d(a).
\end{aligned}
\]
The first branch is unital and tracial; the second preserves a marked corner
and cannot be used to infer the factor \(II_1\).

\section{Cancellation of the Landauer cost in reversible state inversion}

For a complete variable \(X\) and a projection \(Y=q(X)\),
\[
 H(X)=H(Y)+H(X\mid Y).
\]
If \(q\) is bijective, the fiber is a singleton and \(H(X\mid Y)=0\). If \(q\)
is a reset of the uniform TRIT, the discarded information is \(\log3\) and
the minimum cost is \(k_B\Theta\log3\). For the extension
\(\widehat E_N\) to \(L^1\) of a normal trace-preserving conditional expectation
\(E_N\), and under the stated finiteness hypotheses,
\[
 S_\tau(\widehat E_Nh)-S_\tau(h)
 =D_\tau(h\Vert \widehat E_Nh)\ge0.
\]
The two formulas locate the same phenomenon in different domains: the
cost belongs to the contraction of a fiber, not to the reversible evolution
of the complete state.

\section{Gödel–Turing Obstruction to a Uniform Projective Multiplicity Decider}

Let \(T_e\) be the machine with index \(e\). At level \(n\), define
\[
 X_n^e=\{b_n\}\cup
 \{c_n:T_e\text{ did not halt before }n\}.
\]
Each \(X_n^e\) is computable. The global fiber has two branches if \(T_e\) does not
stop and one if it stops. Therefore
\[
 U_e:\quad \#q_\infty^{-1}(v)=1
 \quad\Longleftrightarrow\quad T_e\downarrow.
\]
A uniform decider for \(U_e\) would decide \(\HALT\). A
recursively axiomatizable theory that is sound and complete for every \(U_e\)
would permit their decision by enumerating in parallel proofs of \(U_e\) and of its
negation. This is the core of the \emph{HMT theorem of uniform incompleteness
for projective multiplicities}; its proof uses a
Godel--Turing reduction. The uniform obstruction does not negate particular selectors
constructed on calibrated HMT images.

\begin{controlbox}[title={Scope of the uniform incompleteness control}]
The Gödel--Turing reduction establishes the uniform obstruction for projective multiplicities. It does not negate the particular selectors constructed on calibrated HMT images, nor does it turn a global decision obstruction into indeterminacy of every concrete fiber.
\end{controlbox}
